\documentclass[10pt,a4paper]{amsart}
\usepackage[margin=19mm]{geometry}
\usepackage{amsmath,amssymb,mathtools}
\usepackage[scr=rsfs,cal=euler]{mathalfa}
\usepackage{xfrac}
\usepackage[unicode,hidelinks,bookmarks=false]{hyperref}
\newcommand{\ie}{i.e.\ }
\newcommand{\cf}{cf.\ }
\newcommand{\resp}{respectively\ }
\newcommand{\Subsection}{\subsection}
\providecommand{\nobreakdash}{\nobreak\hbox{-}}
\setlength{\emergencystretch}{2em}
\multlinegap0pt
\usepackage{bm}
\usepackage{bbm}
\usepackage{dsfont}
\usepackage{xspace}
\usepackage{cancel}
\usepackage{mathtools}
\usepackage{amsthm}
\makeatletter
\@ifundefined{Theorem}{\newtheorem{Theorem}{Theorem}}{}
\@ifundefined{Lemma}{\newtheorem{Lemma}[Theorem]{Lemma}}{}
\makeatother
\renewcommand{\phi}{\varphi}
\title[Describing the Numerical Range and $C$-Numerical Range of Ma]{Describing the Numerical Range and $C$-Numerical Range of Matrices via Their Unitary Orbits and the Joukowsky Transform}
\author{Ryan O'Loughlin}
\date{Source: arXiv:2504.05481v4 (CC BY 4.0)}
\begin{document}
\maketitle

\noindent\emph{Source and licence.} Describing the Numerical Range and $C$-Numerical Range of Matrices via Their Unitary Orbits and the Joukowsky Transform. Version arXiv:2504.05481v4 (CC BY 4.0), distributed under the Creative Commons Attribution 4.0 International licence, as recorded in the arXiv licence metadata for this exact version and shown on the abstract page https://arxiv.org/abs/2504.05481v4. Full rights notice: licenses/arxiv-2504.05481-CC-BY-4.0.txt.\par\smallskip

\noindent\textit{Editorial scope.} Counted unit: the Theorem whose source label is \texttt{3by3J} in the article (source0.tex lines 437--490, source label \texttt{3by3J}), delivered together with the article's own complete original proof of it (source0.tex lines 491--711, one \texttt{proof} environment of 221 lines). No mathematics is written, repaired or re-derived: statement and proof are the article's own text line for line. Required context retained uncounted: ellipsesupport (lines 411--416). Every definition, display and label of those lines is retained unchanged. Omitted from this package and not redistributed: 1--410, 417--436, 712--1450. Nothing inside a retained excerpt is omitted or altered: every retained body is delivered exactly as the article's lines are written, so no omission receipt of any kind accompanies this package. Citations: the retained excerpts carry no citation command at all, so no bibliography entry of the article is transcribed and no reference is redistributed. Reference adaptation: none was needed and none was made; every reference inside the delivered unit resolves inside it, so no unresolved reference or citation is produced. Printed slips kept literal: no printed word, symbol, hypothesis or number of the counted statement or of its proof was altered. Layout and class: the article's own class is \texttt{article} with packages including inputenc, amsmath, babel, url, geometry, graphicx, todonotes, hyperref, amssymb, amsthm, authblk, amsfonts, amscd, latexsym; the wrapper is the lane's pinned \texttt{amsart} editorial wrapper at 10pt on A4, which restates the theorem-like environments the retained text needs and adds the article's own macro definitions that the retained text needs (\texttt{\textbackslash phi}), with extra wrapper packages (bm, bbm, dsfont, xspace, cancel, mathtools). The delivered numbering is the unit's own. Assets: no retained span contains a figure environment, a graphics-inclusion command, a PDF-page inclusion, a table or a TikZ picture; no image file is copied into this package, the package declares no asset dependency and no image file is redistributed with it. Licence: version arXiv:2504.05481v4 of this article is distributed under the Creative Commons Attribution 4.0 International licence, as recorded in the arXiv licence metadata for this exact version and shown on the abstract page https://arxiv.org/abs/2504.05481v4. A full rights notice is delivered at licenses/arxiv-2504.05481-CC-BY-4.0.txt. Characters: every retained excerpt is pure ASCII, so the delivered engine, XeLaTeX with its default Latin Modern text font, reports no missing character.
\begin{Lemma}\label{ellipsesupport}
    Let the ellipse $E$ have boundary parametrised by 
    $$
    t \mapsto a \cos (t) + i b \sin (t) , \, \, a,b >0, \, \, t \in [0,2 \pi].
    $$ Then $ h_E(\phi) = \sqrt{a^2\cos^2(\phi) + b^2\sin^2(\phi)}$.
\end{Lemma}

\begin{Theorem}\label{3by3J}
    Let \begin{equation}\label{Athmain}
        A =  \begin{pmatrix}
        0 & a_{12} & a_{13} \\
        a_{21} & 0 & a_{23} \\
        a_{31} & a_{32} & 0
    \end{pmatrix}.
    \end{equation}
\begin{enumerate}
    \item $W(A)$ is given by
\begin{equation}\label{finalJ3by3}
\left\{
  \begin{aligned}
  & s_1c_1 \,\Big(
c_2 \, J_{a_{12},a_{21}}(e^{i\psi_2}) + s_2 \, J_{a_{13},a_{31}}(e^{i\psi_3})
\Big) + s_1^2 \, s_2c_2 \, J_{a_{23},a_{32}}(e^{i(\psi_3-\psi_2)}) ,\\
  &  \text{such that }\,  \theta_1, \theta_2 \in \left[0, \frac{\pi}{2}\right], \, \psi_2, \psi_3 \in \left[0, 2 \pi\right]
  \end{aligned}
\right\},
\end{equation}
where $s_1, s_2, c_1, c_2$ are shorthands for $\sin ( \theta_1), \sin ( \theta_2), \cos( \theta_1), \cos( \theta_2)$ respectively.



\item Let $\begin{pmatrix}
    0 & a_{12} \\
    a_{21} & 0 
\end{pmatrix}$ and $\begin{pmatrix}
    0 & a_{13} \\
    a_{31} & 0 
\end{pmatrix}$ have numerical range boundaries parametrised by 
$$
\begin{aligned}
    &e^{i \alpha}(X_1\cos(t) + iY_1 \sin(t)), \,  \text{where}  \, X_1,Y_1 \geq 0 \, \text{and} \, \alpha \in [0, 2 \pi] \\
    &e^{i \beta}(X_2 \cos(s) + i Y_2 \sin(s)), \,  \text{where}  \, X_2,Y_2 \geq 0 \, \text{and} \, \beta \in [0, 2 \pi]
\end{aligned}
$$
respectively. If $a_{32} = a_{23}=0$, then $W(A)$ is an ellipse with boundary parametrised by

$$
t \mapsto \frac{e^{i \gamma}}{2}\left(\sqrt{\lambda_1} \cos (t) + i \sqrt{\lambda_2}\sin (t)\right)
$$
where
\[
\begin{aligned}
&\gamma = \frac12 \arg\Big( (X_1^2-Y_1^2)e^{i2\alpha} + (X_2^2-Y_2^2)e^{i2\beta} \Big),\\
&\lambda_{1,2} = \frac{X_1^2 + Y_1^2 + X_2^2 + Y_2^2}{2} \; \\
&\pm \frac{1}{2}\sqrt{(X_1^2-Y_1^2)^2 + (X_2^2-Y_2^2)^2 + 2(X_1^2-Y_1^2)(X_2^2-Y_2^2)\cos\big(2(\alpha-\beta)\big)}.
\end{aligned}
\]
(Here $\lambda_1$ is chosen with the $+$ sign and $\lambda_2$ with the $-$ sign.)
\end{enumerate} 

\end{Theorem}

\begin{proof}


\begin{enumerate}
    \item 


For any $h = \begin{pmatrix}
    z_1 \\
    z_2 \\
    z_3
\end{pmatrix}$ of unit norm set $|z_1|=\cos (\theta_1)$. Then $|z_2|^2 + |z_3|^2 = \sin^2(\theta_1)$, so $|z_2|=\sin (\theta_1)\cos (\theta_2)$, $|z_3|=\sin (\theta_1)\sin (\theta_2)$ 
for some $\theta_2 \in[0,\tfrac{\pi}{2}]$. Write $z_j = e^{i\phi_j}|z_j|$ for phases $\phi_j\in[0,2\pi)$, 
so
\[
h=
\begin{pmatrix}
e^{i\phi_1}\cos (\theta_1) \\
e^{i\phi_2}\sin (\theta_1)\cos (\theta_2) \\
e^{i\phi_3}\sin (\theta_1)\sin (\theta_2)
\end{pmatrix}.
\]
Thus every $h$ of unit norm takes this form for some $\phi_1, \phi_2, \phi_3, \in [0, 2 \pi], \ \theta_1, \theta_2 \in [0, \frac{\pi}{2} ]$.


Denoting $$ \tilde{h} = e^{-i\phi_1}
h =  \begin{pmatrix}
\cos (\theta_1) \\
e^{i(\phi_2 - \phi_1)}\sin (\theta_1)\cos (\theta_2) \\
e^{i(\phi_3-\phi_1)}\sin (\theta_1)\sin (\theta_2)
\end{pmatrix},
$$
and observing that $\langle Ah,h \rangle  = \langle A\tilde{h},\tilde{h} \rangle$ we see
\begin{align}
&\left\{ \langle Ah,h\rangle :  \phi_1, \phi_2, \phi_3, \in [0, 2 \pi], \ \theta_1, \theta_2 \in [0, \frac{\pi}{2} ]  \right\} \\
&= \left\{ \langle Ah,h\rangle :  \phi_1=0, \phi_2, \phi_3, \in [0, 2 \pi], \ \theta_1, \theta_2 \in [0, \frac{\pi}{2} ]  \right\}
\end{align}


Furthermore a direct computation shows $\langle A\tilde{h},\tilde{h}\rangle$ is given by
\[
\begin{aligned}
& \cos (\theta_1) \,\sin (\theta_1) \Big(
\cos (\theta_2) \,(a_{12} e^{i\phi_2} + a_{21} e^{-i\phi_2}) 
+ \sin (\theta_2) \,(a_{13} e^{i\phi_3} + a_{31} e^{-i\phi_3}) 
\Big)\\
&+ \sin^2\theta_1 \,\cos (\theta_2) \,\sin (\theta_2) \,(a_{23} e^{i(\phi_3-\phi_2)} + a_{32} e^{i(\phi_2-\phi_3)})\\
&= \cos (\theta_1) \,\sin (\theta_1) \,\Big(
\cos (\theta_2) \, J_{{a_{12}},a_{21}}(e^{i\phi_2}) + \sin (\theta_2) \, J_{a_{13},a_{31}}(e^{i\phi_3})
\Big)
\\
&+ \sin^2\theta_1 \, \cos (\theta_2) \, \sin (\theta_2) \, J_{a_{23},a_{32}}(e^{i(\phi_3-\phi_2)}),
\end{aligned}
\]
from which \eqref{finalJ3by3} follows.


\item Denote
\[
\begin{aligned}
    &E_1^{\text{rot}} = \{ e^{i\alpha} (X_1  \cos (t) + i Y_1  \sin (t)) : t \in [0,2\pi] \}, \\
    &E_2^{\text{rot}} = \{ e^{i\beta} (X_2 \cos (s) + i Y_2 \sin (s)) : s \in [0,2\pi] \},
\end{aligned}
\]
and observe from part $(a)$,
\begin{equation}\label{supportnumrange}
    W(A) = \left\{ \frac{\sin( 2 \theta_1)}{2} \,\Big(
\cos(\theta_2) \, E_{1}^{rot} + \sin(\theta_2) \, E_{2}^{rot})
\Big) :  \ \theta_1, \theta_2 \in \left[0, \frac{\pi}{2}\right] \right\}
\end{equation}
where the summation above is interpreted as a Minkowski summation.
From \eqref{supportnumrange}, we see $h_{W(A)} $ is attained when $\sin (2 \theta_1) = 1$ and furthermore since the support function of the Minkowski sum of two sets is the sum of the support function on the individual sets we have
\[
h_{W(A)}(\phi) = \frac{1}{2} \left(\max_{\theta_2 \in [0, 2 \pi]} \big( \cos(\theta_2) \, h_{E_1^{\text{rot}}}(\phi) + \sin( \theta_2)\, h_{E_2^{\text{rot}}}(\phi) \big) \right).
\]
Viewing the above as the (real) inner product  $$\left\langle \begin{pmatrix}
    h_{E_1^{\text{rot}}}(\phi) \\
    h_{E_2^{\text{rot}}}(\phi)
\end{pmatrix},  \begin{pmatrix}
    \cos( \theta_2) \\
    \sin( \theta_2)
\end{pmatrix} \right\rangle
$$
and applying Cauchy-Schwarz, we obtain
\begin{equation}\label{support2}
h_{W(A)}(\phi)=\frac{1}{2} \left( \sqrt{ h_{E_1^{\text{rot}}}(\phi)^2 + h_{E_2^{\text{rot}}}(\phi)^2 } \right).
\end{equation}





By Lemma \ref{ellipsesupport} for $E_1 = e^{-i \alpha}E_1^{\text{rot}}, \, E_2 = e^{-i \beta}E_2^{\text{rot}}$
\[
h_{E_1}(\phi) = \sqrt{X_1^2 \cos^2 (\phi) + Y_1^2 \sin^2 (\phi)},\quad
h_{E_2}(\phi) = \sqrt{X_2^2 \cos^2 (\phi) + Y_2^2 \sin^2 (\phi)},
\]
so 
\[
\begin{aligned}
&h_{E_1^{\text{rot}}}(\phi) = \sqrt{X_1^2 \cos^2(\phi-\alpha) + Y_1^2 \sin^2(\phi-\alpha)}\\
&h_{E_2^{\text{rot}}}(\phi) = \sqrt{X_2^2 \cos^2(\phi-\beta) + Y_2^2 \sin^2(\phi-\beta)}.
\end{aligned}
\]



For $E_1^{\mathrm{rot}}$ and $E_2^{\mathrm{rot}}$ we write the squared support function in quadratic form.  
Let 
\[
u = \begin{pmatrix} \cos\phi \\  \sin\phi \end{pmatrix}.
\]

Then
\[
h_{E_1^{\mathrm{rot}}}(\phi)^2 
= u^T \, R(\alpha)
\begin{pmatrix} X_1^2 & 0 \\  0 & Y_1^2 \end{pmatrix}
R(\alpha)^T \, u,
\]
\[
h_{E_2^{\mathrm{rot}}}(\phi)^2 
= u^T \, R(\beta)
\begin{pmatrix} X_2^2 & 0 \\  0 & Y_2^2 \end{pmatrix}
R(\beta)^T \, u,
\]
where \[
R(\alpha) =
\begin{pmatrix}
\cos\alpha & -\sin\alpha \\
\sin\alpha & \cos\alpha
\end{pmatrix}
\]
is the unitary rotation matrix.


Hence from \eqref{support2}, $h_{W(A)}(\phi)^2 = \frac{1}{4}u^T M u,$ where
\begin{equation}\label{M}
    M = R(\alpha)\begin{pmatrix}X_1^2&0\\0&Y_1^2\end{pmatrix}R(\alpha)^T
\;+\;R(\beta)\begin{pmatrix}X_2^2&0\\0&Y_2^2\end{pmatrix}R(\beta)^T.
\end{equation}


As $M$ is the sum of two real symmetric positive definite matrices, $M$ is also real symmetric and positive definite, so it can be unitarily diagonalised as
\[
M = R(\gamma)
\begin{pmatrix}\lambda_1 & 0 \\  0 & \lambda_2\end{pmatrix}
R(\gamma)^T, \qquad \lambda_1,\lambda_2 \geq 0.
\]

Therefore
\[
h_{W(A)}(\phi) =\frac{1}{2} \sqrt{\lambda_1 \cos^2(\phi-\gamma) + \lambda_2 \sin^2(\phi-\gamma)}.
\]
This is exactly the support function of an ellipse centered at the origin, rotated by angle $\gamma$, with semi-axes $\frac{\sqrt{\lambda_1}}{2}, \frac{\sqrt{\lambda_2}}{2}$.  


We now compute the values $\gamma, \lambda_1, \lambda_2$. Denote
\[
M = \begin{pmatrix} m_{11} & m_{12} \\[2mm] m_{12} & m_{22} \end{pmatrix},
\]
then directly computing from \eqref{M} gives
\[
\begin{aligned}
m_{11} &= X_1^2 \cos^2\alpha + Y_1^2 \sin^2\alpha + X_2^2 \cos^2\beta + Y_2^2 \sin^2\beta,\\[1mm]
m_{22} &= X_1^2 \sin^2\alpha + Y_1^2 \cos^2\alpha + X_2^2 \sin^2\beta + Y_2^2 \cos^2\beta,\\[1mm]
m_{12} &= (X_1^2 - Y_1^2)\sin\alpha \cos\alpha + (X_2^2 - Y_2^2)\sin\beta \cos\beta.
\end{aligned}
\]
So
\[
\begin{aligned}
m_{11}-m_{22} &= (X_1^2 - Y_1^2)\cos(2\alpha) + (X_2^2 - Y_2^2)\cos(2\beta),\\[1mm]
2 m_{12} &= (X_1^2 - Y_1^2)\sin(2\alpha) + (X_2^2 - Y_2^2)\sin(2\beta),\\[1mm]
m_{11}+m_{22} &= X_1^2 + Y_1^2 + X_2^2 + Y_2^2.
\end{aligned}
\]

As $M$ is symmetric, the eigenvalues \(\lambda_1, \lambda_2\) of \(M\) are

\begin{equation}\label{lambda12}
    \lambda_{1,2} = \frac{m_{11}+m_{22}}{2} \;\pm\; \frac{1}{2}\sqrt{(m_{11}-m_{22})^2 + 4 m_{12}^2}.
\end{equation}

Substituting the $m_{ij}$ values into the above gives
\[
\begin{aligned}
&\lambda_{1,2} = \frac{X_1^2 + Y_1^2 + X_2^2 + Y_2^2}{2} \; \\
&\pm \frac{1}{2}\sqrt{(X_1^2-Y_1^2)^2 + (X_2^2-Y_2^2)^2 + 2(X_1^2-Y_1^2)(X_2^2-Y_2^2)\cos\big(2(\alpha-\beta)\big)}.
\end{aligned}
\]

As $M$ is symmetric, for $
\theta = \tfrac{1}{2}\arctan\!\left(\frac{2m_{12}}{m_{11}-m_{22}}\right)$ (this is often called the principal axis orientation)
the eigenvectors of $M$ are
\[
v_1 = \begin{pmatrix}
\cos\theta \\ 
\sin\theta
\end{pmatrix},
\qquad
v_2 = \begin{pmatrix}
-\sin\theta \\ 
\cos\theta
\end{pmatrix},
\]
with $v_1$ corresponding to the larger eigenvalue (i.e. the eigenvalue with $+$ in \eqref{lambda12}).



Thus
\[
\gamma = \frac{1}{2} \arctan \left( \frac{(X_1^2-Y_1^2)\sin(2\alpha) + (X_2^2-Y_2^2)\sin(2\beta)}{ (X_1^2-Y_1^2)\cos(2\alpha) + (X_2^2-Y_2^2)\cos(2\beta)} \right),
\]
or equivalently
\[
\gamma = \frac12 \arg\Big( (X_1^2-Y_1^2)e^{i2\alpha} + (X_2^2-Y_2^2)e^{i2\beta} \Big).
\]

\end{enumerate}
\end{proof}

\end{document}
